\section{Important Terms And Nomenclature}

In this section, we will define some important terms and nomenclature that will be used throughout the thesis. This is not an exhaustive list, but it should provide a good starting point for understanding the concepts and terminology used in the thesis.

\vspace{1cm}
\textbf{Drone} - Usually by drone we mean an unmanned multicopter or other unmanned flying vehicle. This is not exactly right, this term should be used for any unmanned autonomous vehicle (UAV). In Merriam-Webster dictionary, a drone is defined as:
\textbf{"an uncrewed aircraft or vessel guided by remote control or onboard computers"~\cite{merriam_webster_drone} }
Because in this project we are focused on water operations, by term drone or UAV we will always mean an unmanned surface vehicle / autonomous unmanned ship.



\vspace{0.5cm}
\textbf{IMU} - Inertial Measurement Unit. A device that measures and reports a body's acceleration and angular velocity. In case of this project, by IMU I will always mean 6 degrees of freedom IMU containing both accelerometers and gyroscopes, which measures acceleration in 3 axes and angular velocity in 3 axes. This is the most common type of IMU used in drones and other autonomous vehicles.


\vspace{0.5cm}
\textbf{SITL} - Software in the Loop. A simulation technique used to test and validate control algorithms and systems before deploying them on actual hardware.

\vspace{0.5cm}
\textbf{HITL} - Hardware in the Loop. A simulation technique that integrates real hardware components into the simulation loop, allowing for more realistic testing and validation of control algorithms and systems.


\vspace{0.5cm}
\textbf{Gazebo} - A 3D simulation environment used for testing and validating control algorithms and systems in a realistic virtual environment. Gazebo is often used in conjunction with ROS (Robot Operating System) to create a complete simulation and testing framework for autonomous vehicles and robots.

\newpage
